% ======================================================================
% Diagrama de flujo — generado automáticamente por generar_diagrama_flujo.py
% NO EDITAR MANUALMENTE; regenerar desde el descriptor JSON.
% ======================================================================
\documentclass[11pt,a4paper]{article}

% ── Codificación, fuentes y lenguaje ───────────────────────────────────────────
\usepackage[utf8]{inputenc}
\usepackage[T1]{fontenc}
\usepackage[default,scale=0.95]{sourcesanspro}
\renewcommand{\familydefault}{\sfdefault}
\usepackage[spanish,es-noshorthands,es-tabla]{babel}

% ── Geometría y espaciado ─────────────────────────────────────────────────────
\usepackage[top=2.2cm, bottom=2.5cm, left=2.2cm, right=2.2cm, headheight=15pt]{geometry}
\usepackage{setspace}
\setstretch{1.18}
\usepackage{parskip}

% ── TikZ (simbolos ISO 5807 / ANSI X3.5) ──────────────────────────────────────
\usepackage{tikz}
\usetikzlibrary{babel, arrows.meta, positioning, shapes.geometric,
                shapes.symbols, decorations.pathmorphing, calc}

% ── Tablas y listas ───────────────────────────────────────────────────────────
\usepackage{booktabs}
\usepackage{tabularx}
\usepackage{array}
\usepackage{enumitem}

% ── Colores, cajas e iconos ───────────────────────────────────────────────────
\usepackage[dvipsnames]{xcolor}
\usepackage{tcolorbox}
\tcbuselibrary{skins, breakable}
\usepackage{fontawesome5}

% ── Secciones con barra de color (infografía) ─────────────────────────────────
\usepackage{titlesec}
\titleformat{\section}
  {\normalfont\LARGE\bfseries\color{azulNoche}}
  {\colorbox{cyanNeon}{\color{fondoOscuro}\thesection}}{0.7em}{}
  [\vspace{2pt}\color{cyanNeon}\rule{\linewidth}{1.8pt}]
\titlespacing*{\section}{0pt}{24pt}{10pt}
\titleformat{\subsection}
  {\normalfont\large\bfseries\color{azulNoche}}
  {\textcolor{cyanNeon}{\thesubsection}}{0.7em}{}
  [\vspace{1pt}\color{grisLinea}\rule{\linewidth}{0.6pt}]
\titlespacing*{\subsection}{0pt}{14pt}{6pt}

% ── Cabeceras y pies de página ────────────────────────────────────────────────
\usepackage{fancyhdr}
\usepackage{lastpage}
\pagestyle{fancy}
\fancyhf{}
\renewcommand{\headrulewidth}{0pt}
\renewcommand{\footrulewidth}{0pt}
\fancyfoot[C]{%
  \begin{minipage}{\linewidth}
    \vspace{2pt}
    \color{cyanNeon}\rule{\linewidth}{1.2pt}\\[2pt]
    \centering\color{grisTexto}\footnotesize
    Página \thepage\ de \pageref{LastPage}
  \end{minipage}%
}

% ── Hiperenlaces ──────────────────────────────────────────────────────────────
\usepackage[colorlinks=true, linkcolor=azulNoche, urlcolor=cyanNeon]{hyperref}
\sloppy
\emergencystretch 3em

% ── Paleta de colores ─────────────────────────────────────────────────────────
\definecolor{fondoOscuro}{HTML}{0D1B2A}
\definecolor{verdeTurquesa}{HTML}{004D40}
\definecolor{azulNoche}{HTML}{1B3A6B}
\definecolor{azulMedio}{HTML}{2A5298}
\definecolor{cyanNeon}{HTML}{00C8E0}
\definecolor{verdeSignal}{HTML}{00B887}
\definecolor{naranjaVivo}{HTML}{FF7043}
\definecolor{rojoAlerta}{HTML}{E53935}
\definecolor{grisPapel}{HTML}{F4F6FA}
\definecolor{grisLinea}{HTML}{C8D0E0}
\definecolor{grisTexto}{HTML}{6B7A99}
\definecolor{amarilloNota}{HTML}{FFB300}

% ── Tarjetas ──────────────────────────────────────────────────────────────────
\tcbset{
  cajaContenido/.style={
    enhanced, breakable,
    arc=5pt, outer arc=5pt,
    colback=grisPapel, colframe=grisLinea,
    boxrule=0.8pt,
    fonttitle=\bfseries\small, coltitle=azulNoche,
    top=6pt, bottom=6pt, left=10pt, right=10pt,
  },
}

% ── Macros ────────────────────────────────────────────────────────────────────
\newcommand{\iconotexto}[2]{\textcolor{cyanNeon}{\faIcon{#1}}~#2}
\newcommand{\iconoBanda}{sitemap}
\newcommand{\bandaTitulo}[3][fondoOscuro]{%
  \noindent
  \begin{tcolorbox}[%
    enhanced, arc=4mm, colback=#1!10!white, colframe=#1, boxrule=0pt,
    borderline west={6pt}{0pt}{cyanNeon},
    top=6pt, bottom=6pt, left=8pt, right=8pt,
  ]
    \noindent
    \begin{minipage}[c]{0.10\linewidth}
      \centering\color{cyanNeon}\faIcon{flask}
    \end{minipage}%
    \hspace{8pt}%
    \begin{minipage}[c]{0.88\linewidth}
      {\fontsize{20}{24}\selectfont\bfseries\color{white}#2}\par\vspace{5pt}%
      \textcolor{cyanNeon!80!white}{\small\faIcon{angle-right}~#3}%
    \end{minipage}
    \vspace{4pt}
  \end{tcolorbox}%
  \vspace{8pt}%
}


  % ── Estilos de flujo (ISO 5807 / ANSI X3.5) ─────────────────────────────
  \tikzset{
    flujo/.style={
      >=Stealth,
      font=\footnotesize\sffamily,
    },
    terminador/.style={
      draw=azulNoche, fill=azulNoche!8!white, rounded corners=3pt,
      text width=1.6cm, align=center, inner sep=2mm,
      font=\footnotesize\sffamily,
    },
    proceso/.style={
      draw=azulNoche, fill=cyanNeon!10!white,
      text width=1.6cm, align=center, inner sep=2mm,
      font=\footnotesize\sffamily,
    },
    entradasalida/.style={
      draw=azulNoche, fill=verdeSignal!10!white,
      trapezium, trapezium left angle=75, trapezium right angle=105,
      text width=1.8cm, align=center, inner sep=2mm,
      font=\footnotesize\sffamily,
    },
    decision/.style={
      draw=azulNoche, fill=amarilloNota!20!white,
      diamond, aspect=1.6, text width=1.5cm, align=center, inner sep=1pt,
      font=\footnotesize\sffamily,
    },
    almacenamiento/.style={
      draw=azulNoche, fill=grisPapel,
      cylinder, shape border rotate=90, aspect=0.3,
      text width=1.7cm, align=center, inner sep=2mm,
      font=\footnotesize\sffamily,
    },
    documento/.style={
      draw=azulNoche, fill=azulMedio!8!white,
      text width=1.8cm, align=center, inner sep=2mm,
      font=\footnotesize\sffamily,
    },
    nota/.style={
      draw=grisLinea, dashed, fill=grisPapel,
      text width=1.8cm, align=center, inner sep=2mm,
      font=\scriptsize\sffamily,
    },
    etiqueta/.style={
      font=\scriptsize\sffamily, fill=white,
      fill opacity=0.75, text opacity=1, inner sep=1pt,
    },
  }


% ── Cabecera del documento ────────────────────────────────────────────────────
\fancyhead[L]{\small\color{azulNoche}\textbf{ELECTRONICA}}
\fancyhead[R]{\small\color{grisTexto}flujo\_ocr\_a\_practicas}
\renewcommand{\headrulewidth}{0pt}

% ─────────────────────────────────────────────────────────────────────────────
\begin{document}

\bandaTitulo{Del OCR del PDF a las ocho practicas}{Tres capas, mismo markdown siempre los mismos .tex, y un log sin lineas ! como condicion de exito}

\vspace{4pt}
\begin{center}
\begin{tcolorbox}[cajaContenido, title={\faIcon{map-signs}~Leyenda de símbolos---ISO 5807 / ANSI X3.5}]
\small
Óvalo = \textbf{terminador} (inicio/fin) $\cdot$
Rectángulo = \textbf{proceso} (acción determinista) $\cdot$
Rombo = \textbf{decisión} (ramas Sí/No) $\cdot$
Paralelogramo = \textbf{entrada/salida} (lectura/escritura de archivos) $\cdot$
Cilindro = \textbf{almacenamiento online} (estado, snapshots) $\cdot$
Rectángulo con base ondulada = \textbf{documento} (sesiones, logs) $\cdot$
Rectángulo punteado gris = \textbf{nota explicativa} (sin acción).
Flujo: arriba$\rightarrow$abajo, izquierda$\rightarrow$derecha.

\faIcon{tags}~Cada símbolo lleva una \textbf{etiqueta} (T/P/D/E/A/N + número, según su tipo);
su significado se expande en la tabla \emph{Leyenda de etiquetas} bajo cada diagrama.
\end{tcolorbox}
\end{center}

\vspace{6pt}

\section{\iconotexto{shield-alt}{Validar y generar}}

\begin{center}
\begin{tikzpicture}[flujo, >=Stealth, x=1cm, y=1cm]
  % Fondo decorativo
  \fill[azulNoche!3!white, rounded corners=4pt]
    (-2.3,1.8) rectangle (5.8,-17.7);
  \node[terminador] (T1) at (0.0,0.0) {T1};
  \node[proceso] (P1) at (0.0,-2.2) {P1};
  \node[decision] (D1) at (0.0,-4.4) {D1};
  \node[terminador] (T2) at (3.5,-4.4) {T2};
  \node[proceso] (P2) at (0.0,-6.6) {P2};
  \node[decision] (D2) at (0.0,-8.8) {D2};
  \node[terminador] (T3) at (3.5,-8.8) {T3};
  \node[proceso] (P3) at (0.0,-11.0) {P3};
  \node[proceso] (P4) at (0.0,-13.2) {P4};
  \node[decision] (D3) at (0.0,-15.4) {D3};

  \draw[->] (T1.south) -- (P1.north);
  \draw[->] (P1.south) -- (D1.north);
  \draw[->] (D1.east) -- (T2.west) node[etiqueta, above, pos=0.5]{Si};
  \draw[->] (D1.south) -- (P2.north) node[etiqueta, left, pos=0.5]{No};
  \draw[->] (P2.south) -- (D2.north);
  \draw[->] (D2.east) -- (T3.west) node[etiqueta, above, pos=0.5]{Si};
  \draw[->] (D2.south) -- (P3.north) node[etiqueta, left, pos=0.5]{No};
  \draw[->] (P3.south) -- (P4.north);
  \draw[->] (P4.south) -- (D3.north);
\end{tikzpicture}
\end{center}

\vspace{6pt}
\begin{tcolorbox}[cajaContenido, title={\faIcon{table}~Leyenda de etiquetas}]
\small
\begin{tabularx}{\linewidth}{@{}>{\centering\arraybackslash}p{1.4cm}X@{}}
\toprule
\textbf{Etiqueta} & \textbf{Significado} \\ \midrule
D1 & --check: validar que las 8 cabeceras PRACTICA n (con tilde, sin tilde o en negrita) y sus cinco secciones esten presentes, sin escribir ni compilar. \\
D2 & Falta alguna practica o seccion, o queda un marcador FIGURA/ILEGIBLE sin convertir, o la matematica tiene dolares impares o display desbalanceado? \\
D3 & El log tiene 0 lineas ! y 0 Overfull hbox, y el compilador devolvio success? \\
P1 & Comprobar que el markdown existe y no esta vacio. Si falta, generarlo antes con flujo\_ocr\_pdf.py sobre el escaneo practicas.pdf. \\
P2 & Parsear: limpiar marcas de pagina y numeros sueltos, dividir en 8 practicas y extraer objetivos, teoria, experimental, materiales, procedimiento y actividades por su TITULO, no por su numero irregular. \\
P3 & Renderizar con el preambulo infografico canonico: viñetas a itemize, tablas markdown a tabular, figuras inline a tcolorbox, ecuaciones \$...\$ y \[ ... \] preservadas por \_tm. \\
P4 & Escribir practica\_01.tex a practica\_08.tex y compilar cada una aislada en .build/, leyendo el .log de cada job. \\
T1 & Inicio: ejecutar flujo\_ocr\_a\_practicas.py con --ocr apuntando al markdown de 56 paginas y --salida al directorio de entregables. \\
T2 &  Informe de cobertura y salir sin tocar archivos; es el paso previo seguro antes de gastar compilacion. \\
T3 &  Codigo 3: el informe lista problemas y encontradas. Se corrige la fuente o se relanza el OCR de esa pagina; nunca se inventa contenido. \\
\bottomrule
\end{tabularx}
\end{tcolorbox}

\section{\iconotexto{check-circle}{Validar, publicar y registrar}}

\begin{center}
\begin{tikzpicture}[flujo, >=Stealth, x=1cm, y=1cm]
  % Fondo decorativo
  \fill[azulNoche!3!white, rounded corners=4pt]
    (-2.3,1.8) rectangle (5.8,-15.5);
  \node[terminador] (T4) at (0.0,0.0) {T4};
  \node[decision] (D4) at (0.0,-2.2) {D4};
  \node[terminador] (T5) at (3.5,-2.2) {T5};
  \node[proceso] (P5) at (0.0,-4.4) {P5};
  \node[proceso] (P6) at (0.0,-6.6) {P6};
  \node[almacenamiento] (A1) at (0.0,-8.8) {A1};
  \node[entradasalida] (E1) at (0.0,-11.0) {E1};
  \node[terminador] (T6) at (0.0,-13.2) {T6};

  \draw[->] (T4.south) -- (D4.north);
  \draw[->] (D4.east) -- (T5.west) node[etiqueta, above, pos=0.5]{Si};
  \draw[->] (D4.south) -- (P5.north) node[etiqueta, left, pos=0.5]{No};
  \draw[->] (P5.south) -- (P6.north);
  \draw[->] (P6.south) -- (A1.north);
  \draw[->] (A1.south) -- (E1.north);
  \draw[->] (E1.south) -- (T6.north);
\end{tikzpicture}
\end{center}

\vspace{6pt}
\begin{tcolorbox}[cajaContenido, title={\faIcon{table}~Leyenda de etiquetas}]
\small
\begin{tabularx}{\linewidth}{@{}>{\centering\arraybackslash}p{1.4cm}X@{}}
\toprule
\textbf{Etiqueta} & \textbf{Significado} \\ \midrule
A1 & Guardar la vista por corrida en .tmp/run\_state\_<run\_id>.json con los pasos 1\_generar y 2\_validar, cada uno con exit\_code y timestamp. \\
D4 & Hay algun ! en el log? pdflatex en nonstopmode genera el PDF aunque el documento este roto, asi que la existencia del PDF NO prueba que compilo bien. \\
E1 & Emitir el informe JSON por stdout con status, artefactos y sha\_conjunto. Emitir alerta audible de exito o espera segun el codigo. \\
P5 & Copiar los ocho PDF al directorio de salida junto a sus .tex. Cada artefacto lleva numero, titulo, sha256\_16 y rutas. \\
P6 & Validar la entrega: exige un .tex y un .pdf por practica pedida con mas de 1000 bytes. Calcular sha\_conjunto sobre los ocho documentos. \\
T4 & Entrada a la segunda seccion: el resultado de la decision D3 sobre los logs de compilacion. \\
T5 &  Codigo 4: no se copia el PDF, se informa la linea del log. Causas tipicas documentadas en .agent/latex.md. \\
T6 & Fin: ocho enunciados LaTeX/PDF listos para el curso, reproducibles desde el mismo markdown. \\
\bottomrule
\end{tabularx}
\end{tcolorbox}


\section{\iconotexto{sticky-note}{Notas}}
\begin{itemize}
\item El markdown del OCR es la fuente de verdad del CONTENIDO de las practicas, decision del usuario 2026-11-21: contiene los enunciados que la Universidad exige y sustituye a los .tex derivados de curso.yaml para P1-P8
\item El OCR no es fuente tipografica: trae 34 marcadores FIGURA, una portada ILEGIBLE y tablas comprimidas. El contenido se conserva en cajas de aviso; la imagen real se consulta en el PDF original, nunca se inventa
\item La decision la toma el codigo de execution/generar\_practicas\_desde\_ocr.py, no el chat. El orquestador solo valida entradas, lanza el script, valida salida, guarda la vista por corrida y alerta
\item Cero gasto de LLM y cero red: el OCR ya esta hecho. Aqui solo se parsea markdown y se compila LaTeX en local; no se consume OPENROUTER\_API\_KEY ni cuota de Gemini
\item Un ! en el log es fallo aunque el PDF exista: pdflatex en nonstopmode genera el PDF de un documento roto. Se exigen 0 lineas ! y 0 Overfull hbox antes de copiar el PDF
\item Determinismo: mismo markdown produce el mismo sha\_conjunto. El mismo descriptor SIEMPRE produce el mismo tier en el enrutador; aqui la misma entrada produce los mismos ocho hashes
\end{itemize}

% ─────────────────────────────────────────────────────────────────────────────
\end{document}
